\SourceLeaf{251}%
\Needspace{7\baselineskip}
\begin{center}
\textsc{Gauss} an \textsc{Gerling.} Göttingen, 17. April 1844.
\end{center}

Die Art, wie Sie die Volumengleichheit bloß symmetrischer, nicht zugleich kongruenter Körper beweisen, hat mir viel Vergnügen gemacht. Man könnte den entscheidenden Punkt hervorheben, indem man sagt,

1) dass man auf die Pyramiden aufmerksam macht, deren symmetrische Gegenstücke mit ihnen zugleich kongruent sind. Das sind nämlich diejenigen, bei denen zwei Seitenflächen senkrecht auf einer dritten stehen und in ihren Schnitten mit dieser gleich lange Kanten ergeben, und

2) dass man nachweist, wie sich jede Pyramide in 12 Pyramiden der eben bezeichneten Art zerlegen lässt.

Ob die Beweisart neu ist, kann ich übrigens nicht mit Gewissheit verbürgen. Sehen Sie wenigstens in \textsc{Legendres} Geometrie nach. Dort wird versucht, manches strenger oder einfacher zu beweisen, als es sonst geschehen ist. Ich selbst habe das Buch im Augenblick nicht zur Hand und besitze es auch nicht.

Da jener Satz nun nicht mehr davon betroffen ist, muss ich mein \textls{Bedauern} auf die anderen Sätze der Stereometrie beschränken, die noch von der Exhaustionsmethode abhängen, wie Euklid XII. 5. Vielleicht lässt sich auch hier noch manches verbessern. Im Augenblick habe ich keine Zeit, weiter über den Gegenstand nachzudenken. \mbox{.\ .\ .\ .\ .\ .}

\Needspace{7\baselineskip}
\begin{center}
[\textsc{Gerling} \emph{an} \textsc{Gauss}. \emph{Marburg, 7. Juli 1844.}]
\end{center}

\{\mbox{.\ .\ .\ .\ .\ .\ .} Was zunächst die Symmetrie betrifft, bei der ich im vorletzten Brief nur die »Tatsache« angegeben habe, so scheint mir der \textls{Grund} in unserer Anschauung der Geraden und der Ebene zu liegen. Kraft dieser Anschauung können wir uns in der Ebene an einem Punkt einer Geraden nur \textls{eine} Senkrechte vorstellen, diese aber nach \textls{beiden Seiten} bis ins Unendliche verlängert, und ebenso\space%
